\section{From probability laws to observed data}

The previous chapters studied probability laws as mathematical models. A random
variable \(X\) had a probability law described by a probability mass function, a
density function, or a cumulative distribution function. We defined quantities
such as
\[
\E[X],\qquad \Var(X),\qquad \sigma_X,
\]
and we interpreted them as properties of the probability model.

This chapter changes the point of view. We now begin with data: a finite list of
observed values. The central question is not only what model could have produced
those values, but also how to describe the values that have actually been
observed.

\begin{summarybox}
A probability law is an ideal mathematical object. A dataset is a finite list of
observations. Descriptive statistics are numerical and graphical summaries
computed from that finite list.
\end{summarybox}

This distinction is essential. If a model says that
\[
X\sim N(0,1),
\]
then the model has theoretical expectation \(0\) and theoretical variance \(1\).
But an observed sample generated from that model might be
\[
-0.31,\quad 0.84,\quad 1.12,\quad -1.45,\quad 0.07.
\]
The average of these five observed values will not usually be exactly \(0\), and
their empirical variance will not usually be exactly \(1\). This is not a
contradiction. The model describes the random mechanism; the sample is one finite
outcome of that mechanism.

\subsection*{The same words, but a different status}

Many words from the previous chapters appear again:
\begin{itemize}
    \item distribution,
    \item mean,
    \item variance,
    \item standard deviation,
    \item quantile,
    \item cumulative distribution function.
\end{itemize}

However, their status changes. Earlier, they were properties of a probability
law. Now they are summaries computed from data.

For example, the theoretical expectation
\[
\E[X]
\]
is a property of the random variable \(X\). The sample mean
\[
\bar{x}=\frac1n\sum_{i=1}^n x_i
\]
is a number computed from observed values
\[
x_1,x_2,\ldots,x_n.
\]
These two objects are related, but they are not the same object.

\begin{remarkbox}
A common mistake is to write as if the sample mean \(\bar{x}\) were the same as
\(\E[X]\). The sample mean is computed from data. The expectation belongs to a
probability model.
\end{remarkbox}

\subsection*{Why empirical summaries matter}

Empirical summaries serve several purposes. They help us see the structure of a
dataset, detect unusual observations, compare groups, and form questions about
possible probability models. They also prepare the ground for statistical
inference, where one uses data to reason about an unknown population or an
unknown model parameter.

At this stage, however, we are not yet proving inferential statements. We are
learning the language of data description. The goal is to understand what a
summary measures before using it for estimation or decision making.

\subsection*{A small example}

Suppose the observed waiting times, in minutes, are
\[
2.1,\quad 3.4,\quad 1.8,\quad 4.0,\quad 2.7.
\]
From these five numbers, we can compute a sample mean, a sample variance, a
median, and other summaries. We can also make a histogram or an empirical
cumulative distribution function.

But these five numbers are not themselves an exponential distribution, a normal
distribution, or any other theoretical distribution. They are data. A theoretical
model may later be proposed to explain or approximate how such data arise. The
empirical summaries describe what has been observed.

\subsection*{The guiding comparison}

Throughout this chapter, we repeatedly compare theoretical and empirical objects:
\begin{center}
\begin{tabular}{lll}
\toprule
Theoretical probability & Empirical data & Interpretation \\
\midrule
Probability law of \(X\) & empirical distribution & how mass is assigned \\
CDF \(F_X\) & empirical CDF \(F_n\) & accumulated probability or proportion \\
\(\E[X]\) & sample mean \(\bar{x}\) & center \\
\(\Var(X)\) & sample variance & spread \\
Theoretical quantile & empirical quantile & probability position \\
Density & histogram estimate & shape of distribution \\
\bottomrule
\end{tabular}
\end{center}

The table is only a guide. It does not say that the empirical object is equal to
the theoretical object. It says that empirical summaries are data-based
counterparts of the theoretical concepts developed earlier.

\begin{theorembox}
The main conceptual task of this chapter is to keep two levels separate while
understanding their relationship: the level of the probability model and the
level of the observed sample.
\end{theorembox}
